\section{Conclusion}
\label{sec:conclusion}

Over the constrained domain $z\le1.8$, the two Planck-compatible
references are much closer to each other than
either is to~R: across the constrained bins the C1--C2 RMS difference
is about 1--8\% of either R--CMB difference. Changing between these two Planck-compatible references
changes the inferred difference only slightly (reference sensitivity,
not independent confirmation).

The two R--CMB difference fields over $z\le1.8$ are nearly identical:
they share one dominant direction
(cosine~\GeometryCosineSimilarity{}, principal angle
\GeometryPrincipalAngle$^{\circ}$) with cross-branch transfer
degradation fraction \GeometryTransferDegradation{} or less, reaching
largest amplitude near $z\sim\GeometryPeakZ{}$ without that redshift
being claimed as a fundamental physical scale.

A five-coordinate $m3$ basis provides a compact empirical
description over $z\le1.8$ ($\chi^{2}=\BridgeChiTwo{}$, local
projected Jacobian rank~5 of~5 fitted coordinates),
but its selection is not itself statistically significant under the
null bootstrap
($P(\mathrm{select}\ge m3)=\CvConstrBootProb{}$), which tests the complexity of
the representation rather than the existence of the numerical
difference between frozen endpoint histories. Continuation to
Ly$\alpha$ at $z=2.33$ fails for the compact extrapolator, bounding
the validated domain to $z\le1.8$.

The physical origin remains open: no dynamics, no metric theory,
and no tension resolution are established here. R is a
phenomenological low-redshift endpoint, not presented as a complete
cosmological model.

The resulting target is the low-redshift distance-sector behavior
represented by R over $z\le1.8$, together with the requirement that
any physical completion preserve the observed CMB acoustic structure;
R itself is not asserted to be such a completion.
